\section{Extended related work}
\label{app:related}

This appendix records background that informed the design but that the main
argument does not depend on. Nothing here is used to support a result in
\S\ref{sec:results}; readers who want only the load-bearing literature should
read \S\ref{sec:related} and stop there.

\subsection{Further informal sources on the game}

Beyond Pagat \cite{pagat} and Develin \cite{develin}, the descriptive corpus
is thin. Pagat's tactics list covers deliberate misses that inform a partner,
locking out a dangerous opponent by never asking them anything, and exhausting
one rank across all opponents before switching rank; it also records a
partner-signalling convention and a ``Forced Claims'' rule, both attributed to
named contributors. Develin's chapter contains the corpus's only quantitative
endgame discussion, in which the alternative to a coin-flip declaration would
need a success probability greater than one. The Wikipedia entry
\cite{wikiliterature} names the ``stalemate breaker'', and secondary strategy
pages \cite{dorsa,strategy-site} restate blackballing, asking the asker, and
an endgame delegation heuristic. We use these as evidence about how the game
is played and talked about, not as sources of quantitative claims. The two
substantial sources also disagree with each other on this point: Develin's
chapter forbids pre-agreed conventions outright, where Pagat records a
partner-signalling convention as an accepted tactic.

On Somani's agent \cite{somani,somani-server}: it is a Python implementation
with a complete rules engine and a sound deduction fixed point, and the
four-player, 48-card variant it implements makes the declaration problem
structurally easier than the one studied here. With a single teammate, naming
the allocation of a six-card half-suit ranges over $2^6$ possibilities and is
frequently forced by the deduction fixed point; with two teammates it ranges
over $3^6 = 729$, and the allocation question stops being a formality. The
published learner regresses a hand-written per-outcome score rather than a
terminal game result, which penalises the deliberate miss that both Pagat and
Develin recommend, and its move generator excludes known misses except as a
fallback. Reading the published source, the terminal comparison that would
supply the win/loss signal tests an integer team identity against an enumeration
member, so it never succeeds; on that reading every move in every training game
receives the losing reward and the released model never saw a win/loss gradient,
which would make it a greedy maximum-likelihood asker rather than a learner of
the game's objective. We record this as a reading of the source rather than as a
reproduction: we did not run the training. No strength numbers have been
published for it, so it is not usable as a baseline here in any case.

One methodological caveat from the local-best-response literature bears directly
on the probe of \S\ref{sec:lbr} and is recorded here because the probe does not
address it. Lis\'y and Bowling report that restricting the point at which the
response is allowed to deviate changes the measured exploitability substantially,
so a single entry point understates it and the entry point has to be swept
\cite{lisy-lbr}. Our response is fitted over the whole game with no sweep over
where it may begin deviating, which is one of the reasons \S\ref{sec:lbr}'s
figure is a lower bound rather than a bound.

\subsection{Equilibrium computation}

The equilibrium line runs from counterfactual regret minimisation
\cite{zinkevich-cfr} through outcome- and external-sampling variants
\cite{lanctot-mccfr}, CFR+ \cite{tammelin-cfrplus} and discounted CFR
\cite{brown-dcfr} to function approximation
\cite{brown-deepcfr,steinberger-dream}. Two properties of this family bear on
Fish. CFR+ performs relatively poorly in games where some actions are very
costly mistakes \cite{brown-dcfr}, and an incorrect declaration is exactly
such an action; and DREAM's exploration term grows with depth
\cite{steinberger-dream}, while a game in our dialect runs to
\numEventsPerGame{} public events on average against \vthree{} in E3. Student
of Games extends the belief-conditioned search line and identifies Scotland
Yard as the closest published public-observation domain \cite{schmid-sog}. We
did not attempt any of this.

\subsection{Team-game complexity results}

\S\ref{sec:related} names the target concept for a Fish team and states why we
do not compute it; this subsection records the machinery that would be needed
if one did. Celli and Gatti establish the complexity of the team-maxmin
equilibrium with correlation and the worst-case cost of forgoing correlation
\cite{celli-gatti}. Farina and coauthors connect ex-ante team collusion to
extensive-form correlation \cite{farina-collusion}; the team belief DAG and
the two-player conversion give a zero-sum game whose equilibria are
realization-equivalent to a team-maxmin equilibrium with correlation
\cite{zhang-tbdag,carminati-tpi}, at a representation cost exponential in how
much private information distinguishes teammates inside a public state.
Reductions to a continuous-state game over public beliefs \cite{brown-rebel},
building on the common-information reduction of decentralised control
\cite{nayyar-common}, are the other standard route. Learning-based attacks on
team games exist \cite{mcaleer-teampsro,xu-fxp}; the fictitious-cross-play
analysis proves that self-play converges to \emph{local} team equilibria
\cite{xu-fxp}, which is the formal version of the informal complaint that a
self-played agent's asks stop carrying meaning outside its own training
population. Together with the information-set size quoted in
\S\ref{sec:related}, these results are why this report makes no equilibrium
claim of any kind.

\subsection{Cooperative play through public actions}

Hanabi is the reference benchmark for coordination through observable actions
\cite{bard-hanabi}, and is relevant because Fish likewise has no side channel.
The Bayesian Action Decoder made the public belief an explicit object of the
policy \cite{foerster-bad}; the Simplified Action Decoder showed that letting
teammates condition on the greedy rather than the exploratory action is a
prerequisite for conventions to survive $\varepsilon$-greedy training
\cite{hu-sad}; Other-Play demonstrated how much of a self-play score is
convention rather than skill, with the same agents scoring far worse when
paired across independent training runs \cite{hu-otherplay}; and Off-Belief
Learning provides an explicit dial on convention depth together with a
uniqueness result that makes independent runs interoperable \cite{hu-obl}.
Test-time search over a fixed blueprint provides a further improvement
\cite{lerer-sparta}. Bridge bidding is the other model system, an auction that
is literally a learned language: bidding networks improved on a strong
commercial program \cite{rong-bidding}, joint policy search --- which scores
simultaneous changes at several information sets, which is what inventing a
convention requires --- improved a system further \cite{tian-jps}, and
training with human partners improved results again \cite{lockhart-bidding}.

The asymmetry that separates Fish from both is the audience. Hanabi has no
adversary, so every bit of signal is pure profit, and in bridge the harm done
by a leaked auction is diffuse. In Fish a leaked card location is directly
executable: once it is public that a player holds a card, any opponent holding
another card of that half-suit can take it and keep the turn. The Hanabi
construction that most clearly fails to port is the modular hat-guessing code,
which works because each receiver sees every other receiver's hand and can
subtract it \cite{cox-hanabi}; in Fish no player sees any hand. Any future
claim that an agent has learned a convention would also have to survive the
scramble ablation of \cite{lowe-pitfalls}, in which high speaker--listener
consistency coexisted with no causal effect.

\subsection{An informal secrecy argument}

The following is an argument, not a theorem; we state it because it is the
design rationale for building the agent with no private channel at all, and we
have not proved it. Teammate and opponent hands are exchangeable given the
public history. An absolute codebook --- one whose meaning is fixed by public
information alone --- is therefore decoded identically by all five listeners,
so with two teammates against three opponents the Wyner-style secrecy rate of
such a convention is non-positive \cite{wyner,csiszar-korner}. Only a
receiver-relative codebook, whose semantics depend on the receiver's own
holdings, would have positive rate, and it pays for that with ambiguity at the
encoder. The justification here is structural and informal: we give no proof,
and none is claimed. Public-signalling economics predicts a similar outcome,
with cheap talk separating before two audiences only when it is credible
against both and otherwise collapsing toward pooling \cite{farrell-gibbons}.

\subsection{Remaining determinization detail}

\S\ref{sec:related} states that averaging perfect-information values over deals
sampled from the public history collapses to a hit-probability heuristic in this
game; the collapse is the following. Almost every reachable position carries the
same perfect-information value, because the team on turn takes everything it can
reach, so with $t$ ranging over opponents and $c$ over askable cards and $h$ the
public history,
\[
  \arg\max_{(t,c)} \; \mathbb{E}\bigl[V^{\mathrm{PI}}(t,c)\bigr]
  \;\approx\;
  \arg\max_{(t,c)} \; \Pr[\,t \text{ holds } c \mid h\,].
\]
The right-hand side is one ply deep and blind to the information an ask leaks,
the information a refusal supplies and the option value of postponing a
declaration, which is why the hit probability enters \S\ref{sec:policy} as one
feature among several rather than as an engine.

Determinization for bridge was proposed well before it was made to work
\cite{levy}, and constrained deal generation in practice is done by rejection:
deal against suit-length restrictions, then discard deals whose auction the
program's own bidding module would not have produced
\cite{ginsberg-gib,pavlicek}. Frank and Basin named the two mechanisms that
make the method unsound however many worlds are drawn: \emph{strategy fusion},
in which the searcher plays a different strategy in each sampled world
although a real player must commit to one strategy per information set, and
\emph{non-locality}, in which a node's value depends on parts of the tree
outside its own subtree \cite{frank-basin-ai}. Long, Sturtevant, Buro and
Furtak characterise when determinization nevertheless performs well, through
leaf correlation, bias and a disambiguation factor \cite{long-pimc}, and
Information Set Monte Carlo Tree Search replaces the independent trees with a
single tree over information sets, redeterminizing each iteration
\cite{cowling-ismcts}. The repairs that work are those attacking non-locality
as well as strategy fusion \cite{frank-basin-aaai}. Information Set Monte
Carlo Tree Search has been extended by re-determinizing at non-root seats and
by self-determinization, which allows a form of bluffing to emerge
\cite{goodman-ris,cowling-bluffing}; it is nevertheless unsound, and its
measured exploitability can \emph{increase} with additional computation
\cite{lisy-oos}. A Spades post-mortem traced blunders to a sampler that placed
a card of true probability $1/3$ into $87\%$ of its sampled worlds
\cite{whitehouse-spades}, which is the failure mode our oracle test (E15) is
designed to detect. One smaller result survives into our design: Go Fish sits
at the maximum of Zuckerman, Felner and Kraus's Opponent Impact measure
\cite{zuckerman-mpmix}, so the choice of which opponent to ask is a
first-class decision rather than a tie-break, and our ask features treat it as
one.

\subsection{Sampling and history filtering}

The matrix-scaling approximation \S\ref{sec:related} adopts for the Fast path
carries a guarantee: the deterministic strongly polynomial analysis of Linial,
Samorodnitsky and Wigderson bounds the permanent within a factor $e^n$
\cite{lsw}. Of the permanent methods \S\ref{sec:related} sets aside, Ryser's
inclusion--exclusion formula is exact in $O(2^n n)$ time \cite{ryser}, which
is unusable at $n = 54$, and Glynn's sign-vector formula is the other exact
method \cite{glynn}; the Jerrum--Sinclair--Vigoda chain is a fully polynomial
randomised approximation scheme \cite{jsv} whose constants make it impractical
at this scale, and a practical study measured it at tens of thousands of
seconds on instances where Ryser's formula takes milliseconds
\cite{newman-vardi}. Sequential importance sampling is the standard
contingency-table alternative \cite{chen-sis}, with the documented warning
that it can underestimate by an exponential factor under any row or column
ordering while appearing to have converged \cite{bezakova-sis}. This is one
reason \S\ref{sec:belief} uses a direct construction rather than an
importance-weighted one. More generally, constructing a history consistent
with a public state is FNP-complete in the worst case, and enumeration is
polynomial only for sparse public trees \cite{solinas-history}; Fish escapes
this because its public record determines every card movement, so the only
object to be reconstructed is the initial deal.

\subsection{Variance reduction, and why we do not use it}

Poker supplies a family of control variates: advantage-sum estimators
\cite{zinkevich-divat}, importance sampling over imaginary observations
\cite{bowling-is}, MIVAT's least-squares value function \cite{white-mivat},
and AIVAT \cite{burch-aivat,avaivat}. The reported gains are large in
two-player settings and smaller in six-player ones, which is the closer
analogue to our setting. We deliberately do not use AIVAT or MIVAT. AIVAT is
unbiased for \emph{any} value function, which is a constraint on its
expectation rather than on any realised estimate; \cite{kim-sandholm} tuned
such a heuristic on a published six-player dataset until all fourteen players
simultaneously appeared to be large winners, which the zero-sum structure
makes impossible. A learned control variate fitted by us, on our own agent, on
our own evaluation data, is the exact pathology this literature warns about,
and the duplicate-block design of \S\ref{sec:protocol} already removes the
largest part of deal variance without introducing a fitted estimator. That
design is the standard remedy rather than a novel one: duplicate bridge
replays each board with the same cards in the same seats, and the computer
poker competition added common seeds across pairings for the same reason
\cite{acpc}.

\subsection{Rating systems and sequential testing}

Bradley--Terry-style rating models \cite{coulom-whr,herbrich-trueskill} have
two documented failure modes when the rated population is a training run's own
checkpoints: they are not redundancy-invariant, so duplicating one agent
redistributes every rating, and they cannot represent cyclic dominance.
Maximum-entropy Nash averaging over the antisymmetric logit matrix
\cite{balduzzi-reeval} and empirical-game response graphs
\cite{omidshafiei-alpharank} are the proposed alternatives. We report
Bradley--Terry ratings in \S\ref{sec:results-pop} as a compact summary of a
fixed round-robin, not as a claim about a strategy space.

Two further traps shaped \S\ref{sec:protocol}. Monitoring an interval computed
for a fixed sample size and stopping when it looks favourable is not a test;
in simulation this produced a false-positive rate far above nominal under the
null \cite{avaivat}, which is why sequential designs on normalised Elo exist
\cite{nelo}. Our seed banks are fixed in advance and reported at their planned
size. Finally, refining an abstraction does not monotonically reduce
exploitability \cite{waugh-abstraction}, which is the general form of the
observation in \S\ref{sec:ablations} that substituting a more accurate belief
into a policy fitted around a less accurate one does not improve it.
